Fix the same deterministic code pair \((\widehat g_n,\widehat q_n)\), its clipped pair \((\bar g_n,\bar q_n)\), and the same nested radii \((\varepsilon_{1,n},\varepsilon_{2,n})\) as external inputs to the two separate branches. Along every sequence of indices \(n\to\infty\) for which \(n\ge2\), \(\mathcal P_{\mathrm{NG},n}\neq\varnothing\), and \(\varepsilon_{1,n}\le\varepsilon_0\), the broad separated non-Gaussian branch satisfies
\[
 \mathfrak R_n^{\mathrm{NG}}\asymp n^{-1}.
\]
This includes fixed nonzero \(\varepsilon_{1,n}\); the witnessing statistic \(\widehat\theta_{\mathrm{spec},n}\), as well as membership in \(\mathcal P_{\mathrm{NG},n}\), imposes no restriction on and makes no use of \(\widehat q_n\), \(\bar q_n\), or \(\varepsilon_{2,n}\). Separately, every \(P\in\mathcal P_{\mathrm G,n}^{\mathrm{JMS}}\) in the simultaneous bounded-observed-outcome Gaussian intersection defined in \(\text{def:gaussian-class}\) satisfies \(\theta_0(P)=0\). Whenever that intersection is nonempty,
\[
 \mathfrak R_n^{\mathrm G}
 =\mathfrak M_{n,1-\gamma}^{\mathrm G}
 =0.
\]
Thus \(\text{prop:bounded-outcome-gaussian-degeneracy}\) diagnoses an incompatibility internal to the simultaneous exact-PLM, bounded-\(Y\), and nondegenerate conditional-Gaussian assumptions stated in Jin--Mackey--Syrgkanis Theorem 3.2. The positive lower-bound claim from that theorem is not relied upon here. No equality or inclusion between the two branches is asserted.